\appendix

% ============================================================
\section{Lean~4 Proof Artifacts}
\label{app:lean}
% ============================================================

All formal results used in the paper are verified in Lean~4 using Mathlib.
The development contains five modules: \texttt{Basics.lean}, \texttt{Regret.lean}, \texttt{JensenGap.lean}, \texttt{FacetActivation.lean}, and the umbrella import \texttt{ResidualExpertiseNotDecisionValue.lean}.
Table~\ref{tab:lean-map} maps the paper's claims to Lean declarations.

\begin{table}[ht]
\centering
\caption{Map from paper statements to Lean~4 declarations.}
\label{tab:lean-map}
\small
\begin{tabular}{@{}llll@{}}
\toprule
Paper & Type & Lean declaration & Module \\
\midrule
$V(b) = \max_a r_a \cdot b$ & Def & \texttt{terminalValue} & Regret \\
$\operatorname{regret}(a,b) \geq 0$ & Lem & \texttt{regret\_nonneg} & Regret \\
$\operatorname{regret}=0 \Leftrightarrow a$ optimal & Lem & \texttt{regret\_eq\_zero\_iff} & Regret \\
$\E[b_{\mathrm{post}}]=b$ & Lem & \texttt{prior\_eq\_weighted\_posterior} & Basics \\
Jensen gap $\geq 0$ & Thm & \texttt{jensenGap\_nonneg} & JensenGap \\
Thm~\ref{thm:voi-br} (VoI $=$ $\E[\BR]$) & Thm & \texttt{jensenGap\_eq\_weighted\_regret} & FacetActivation \\
Thm~\ref{thm:facet-bridge} & Thm & \texttt{facet\_activation\_bridge} & FacetActivation \\
\bottomrule
\end{tabular}
\end{table}

\paragraph{Primitives.}
\texttt{Basics.lean} defines \texttt{Belief}, \texttt{ObsKernel}, \texttt{marginalProb}, and \texttt{posterior}, and proves the posterior martingale identity \texttt{prior\_eq\_weighted\_posterior}.
\texttt{Regret.lean} defines \texttt{terminalValue} and \texttt{regret} and proves non-negativity and the zero-regret optimality characterization.
\texttt{JensenGap.lean} defines Jensen gaps over posterior lotteries and proves non-negativity for belief-convex functions.

\paragraph{Paper Theorems.}
\texttt{FacetActivation.lean} contains the two formal claims used in the text.
The weighted-regret identity \texttt{jensenGap\_eq\_weighted\_regret} rewrites the value of information as expected regret of retaining the prior-optimal action after the posterior update.
The bridge theorem \texttt{facet\_activation\_bridge} states that this expected regret is strictly positive exactly when some reachable posterior strictly prefers a rival action.

% ============================================================
\section{Synthetic Validation: Per-Configuration Results}
\label{app:exp1-detail}
% ============================================================

Table~\ref{tab:exp1-detail} reports per-configuration results across the three estimator-validation conditions of Section~\ref{sec:exp1}.
False-positive rate is zero in every cell: across all 18 configurations the plug-in estimator never assigned positive boundary regret to an instance with $\BR(m,h)=0$.
Under the positive control C1, $\rho(\BRhat,\BR)$ degrades smoothly as the action/state space grows ($G$ from $2$ to $5$, $K$ from $2$ to $10$, giving $0.969 \to 0.745$), reflecting harder channel estimation, not a failure of the identity.
Under the strong null C3, the human signal is by construction strongly predictive ($\Delta\mathrm{LL}$ up to $0.738$), yet the estimator continues to assign zero boundary regret to every instance, the geometric prediction of Theorem~\ref{thm:facet-bridge}.

\begin{table}[htbp]
\centering
\caption{Per-configuration results for synthetic estimator validation. $G$: decision groups; $K$: latent states; $|A|$: actions in the reward matrix. FPR is zero in every cell and is omitted; $\rho$ is omitted when true boundary regret is constant zero (C2, C3).}
\label{tab:exp1-detail}
\small
\begin{tabular}{@{}lrrrrrr@{}}
\toprule
Condition. & $G$ & $K$ & $|A|$ & $\rho(\BRhat,\BR)$ & $\%\BR{>}0$ & $\Delta$LL \\
\midrule
C1 & 2 & 2  & 2 & 0.969      & 35.5 & 0.226 \\
C1 & 2 & 4  & 2 & 0.984      & 36.4 & 0.230 \\
C1 & 3 & 3  & 3 & 0.941      & 37.9 & 0.403 \\
C1 & 3 & 6  & 3 & 0.903      & 38.2 & 0.420 \\
C1 & 5 & 5  & 5 & 0.862      & 39.8 & 0.577 \\
C1 & 5 & 10 & 5 & 0.745      & 39.8 & 0.592 \\
\midrule
C2 & 2 & 2  & 2 & \text{n/a} &  0.0 & 0.006 \\
C2 & 2 & 4  & 2 & \text{n/a} &  0.0 & 0.006 \\
C2 & 3 & 3  & 3 & \text{n/a} &  0.0 & 0.015 \\
C2 & 3 & 6  & 3 & \text{n/a} &  0.0 & 0.019 \\
C2 & 5 & 5  & 5 & \text{n/a} &  0.0 & 0.035 \\
C2 & 5 & 10 & 5 & \text{n/a} &  0.0 & 0.033 \\
\midrule
C3 & 2 & 2  & 3 & \text{n/a} &  0.0 & 0.447 \\
C3 & 2 & 4  & 3 & \text{n/a} &  0.0 & 0.458 \\
C3 & 3 & 3  & 4 & \text{n/a} &  0.0 & 0.592 \\
C3 & 3 & 6  & 4 & \text{n/a} &  0.0 & 0.613 \\
C3 & 5 & 5  & 6 & \text{n/a} &  0.0 & 0.723 \\
C3 & 5 & 10 & 6 & \text{n/a} &  0.0 & 0.738 \\
\bottomrule
\end{tabular}
\end{table}

\subsection{Sensitivity to Systematic Belief Miscalibration}
\label{app:exp1-miscalibration}

The zero FPRs above concern the unshifted cross-fitted estimates.  To isolate the plug-in estimator's sensitivity to systematic belief error, we use the canonical binary cells ($G=K=2$) of C1 and C3.  C1 mixes positive- and zero-regret cases; C3 is the certified strong null with every case in the FPR denominator.  This binary setting also matches the one-vs-rest downstream task without requiring an arbitrary multiclass direction of perturbation.

For a binary fitted belief $b=(b_0,b_1)$, we apply the signed log-odds offset
\[
T_\delta(b)_1 = \frac{e^\delta b_1}{b_0 + e^\delta b_1},
\qquad
T_\delta(b)_0 = \frac{b_0}{b_0 + e^\delta b_1},
\qquad \delta\in\{-1,+1\}.
\]
Thus $\delta=+1$ multiplies the odds of state 1 by $e$, while $\delta=-1$ divides them by $e$.  We shift $\hat b_x$ and $\hat b_{x,h}$ separately, retain the same simulated cases, fitting protocol, oracle boundary regret, and reward tolerance, and recompute only $\BRhat$ and FPR.  The unshifted baseline is $\delta=0$ and is already reported in the main experiment.

\begin{table}[htbp]
\centering
\caption{Paired binary log-odds stress test of plug-in FPR. Each entry is the percentage of oracle-zero cases assigned $\BRhat>0$ after shifting exactly one fitted belief. $n_0$ is the number of oracle-zero cases, using the same tolerance as FPR.}
\label{tab:exp1-miscalibration}
\small
\begin{tabular}{@{}lrrrrrr@{}}
\toprule
Condition & Baseline & $\hat b_x,-1$ & $\hat b_x,+1$ & $\hat b_{x,h},-1$ & $\hat b_{x,h},+1$ & $n_0$ \\
\midrule
C1 & 0.0 & 50.9 & 49.1 & 0.0 & 0.0 & 6,453 \\
C3 & 0.0 & 0.0 & 0.0 & 40.4 & 23.0 & 10,000 \\
\bottomrule
\end{tabular}
\end{table}

The effect is component- and geometry-dependent.  In C1, perturbing the model-only belief changes the incumbent action and creates false positives among its oracle-zero cases, whereas the same posterior-only shifts do not.  In C3, the safe action makes the model-only shifts inert, but perturbing the human-augmented belief can falsely cross the safe-action facet.  These one-at-a-time shifts intentionally need not preserve the Bayesian coherence relation between $b_x$ and $b_{x,h}$: they are an input-level stress test of separately misspecified deployed estimates, not a perturbation of the exact identity.

This is one interpretable intercept-type stress magnitude, not a general robustness guarantee or an estimate of CheXpert calibration drift.  Calibration fitted and evaluated on genuinely independent (or nested) data could reduce correctable systematic bias, but lower calibration error alone does not guarantee lower boundary-regret FPR near a reward facet.  Standard conformal prediction instead targets coverage of prediction sets; it neither calibrates these point posteriors nor directly controls $\Pr(\BRhat>0\mid\BR=0)$.

% ============================================================
\section{CheXpert Protocol}
\label{app:exp2-protocol}
% ============================================================

\paragraph{Reader Study and Label Construction.}
The five board-certified radiologists (\texttt{bc1}, \texttt{bc2}, \texttt{bc3}, \texttt{bc5}, \texttt{bc7}) provided binary disease-presence ratings under the official CheXpert five-condition protocol (Atelectasis, Cardiomegaly, Consolidation, Edema, Pleural Effusion) on the 500-image frontal-view validation set.
For each (condition, ground-truth reader) pair, the held-out label $y$ is computed by \emph{leave-one-out majority vote}: for reader $r_i$, $y = \Ind[\sum_{j \neq i} v_j \geq 3]$, where $v_j \in \{0,1\}$ are the remaining four ground-truth readers' votes.
This construction enforces $y \perp h$ structurally: the test reader is never a voter in their own ground truth.
The three benchmark readers (\texttt{bc4}, \texttt{bc6}, \texttt{bc8}) were not voters in the original study, so their pairs use all five ground-truth votes (equivalent to the official majority).

\paragraph{Model Belief and Calibration.}
Model probabilities $b_x^{\mathrm{raw}}$ are produced by the \texttt{torchxrayvision} DenseNet-121 \citep{Huang2017,Cohen2022} with the standard X-ray preprocessing pipeline (8-bit normalisation, centre crop, $224{\times}224$ resize).
The primary experiment uses the in-domain checkpoint \texttt{densenet121\allowbreak-res224\allowbreak-chex}; the robustness analyses below use all six publicly available checkpoints.
For each (condition, reader) pair we apply per-pair logit-space temperature scaling: a stratified $30\%$ calibration split (seed $42$), with scalar $T$ found by NLL minimisation in $[0.1,10]$, gives the calibrated belief $b_x = \sigma(\operatorname{logit}(b_x^{\mathrm{raw}})/T)$.

\paragraph{Human-Augmented Posterior.}
On the remaining $70\%$ evaluation split, we estimate $\bpost_{x,h} \approx P(Y=1\mid x,h)$ by $5$-fold cross-fitted logistic regression on features $[\operatorname{logit}(b_x),\,h,\,h\cdot\operatorname{logit}(b_x)]$, with folds stratified on $y$ (\texttt{StratifiedKFold}).
The leave-one-out label construction makes $y \perp h$ structural; cross-fitting additionally makes the augmented-belief estimator independent of the held-out label within each fold, so plug-in $\BRhat$ is not contaminated by within-fold overfitting.

\paragraph{Evaluation.}
For each pair, all per-instance quantities ($\BRhat$, log-loss gain, facet distance) are computed on the held-out fold of each instance and aggregated across folds.
Pair-level summaries (Spearman $\rho$ between log-loss gain and $\BRhat$, top-decile zero-mass fraction) are then aggregated across the $25$ pairs.
Pairs with $\BRhat\equiv 0$ across the entire evaluation set are excluded from the $\rho$ computation; we report the count of valid pairs $n_{\mathrm{v}}$ alongside.

% ============================================================
\section{Reward-Conditional Dissociation: Robustness}
\label{app:exp2-robustness}
% ============================================================

Table~\ref{tab:robustness} extends the dissociation analysis of Section~\ref{sec:exp2} to all six publicly available DenseNet checkpoints and all eight CheXpert radiologists ($6 \times 8 \times 5 = 240$ reader--task pairs).
The ordering R1${\gg}$R2${\geq}$R3 holds on every checkpoint individually: R1 is positive on every checkpoint ($+0.373$ to $+0.530$), R3 ranges from clearly negative (in-domain, $-0.117$ to $-0.040$) to essentially zero (off-domain NIH, $+0.007$), and R1 strictly dominates R2 throughout, with the gap compressing on off-domain pretraining.
Off-domain pretraining (NIH, PadChest) attenuates the R1 signal and inflates R2, consistent with calibration drift moving belief mass toward the lower R2 facet at $1/6$.
The direction of the dissociation, not its absolute magnitude, is the theory-relevant prediction.

\begin{table}[htbp]
\centering
\caption{Reward-conditional dissociation across $6$ DenseNet checkpoints $\times$ $8$ readers $\times$ $5$ conditions $=$ $240$ reader--task pairs.
  Entries report $\bar\rho_{\mathrm{LL}}$, the mean per-pair Spearman correlation between log-loss gain and $\BRhat$, under each reward (R1: thr.\ $1/2$; R2: thr.\ $1/6$; R3: thr.\ $5/6$).
  Bold marks the aggregate dissociation pattern across all $240$ pairs.
}
\label{tab:robustness}
\small
\begin{tabular}{@{}lrrr@{}}
\toprule
Checkpoint & R1 $\bar\rho_{\mathrm{LL}}$ & R2 $\bar\rho_{\mathrm{LL}}$ & R3 $\bar\rho_{\mathrm{LL}}$ \\
\midrule
\multicolumn{4}{l}{\textit{In-domain pretraining}} \\
\texttt{densenet121-res224-chex}      & $+0.489$ & $-0.008$ & $-0.117$ \\
\texttt{densenet121-res224-mimic\_ch} & $+0.496$ & $+0.197$ & $-0.040$ \\
\texttt{densenet121-res224-mimic\_nb} & $+0.495$ & $+0.210$ & $-0.042$ \\
\texttt{densenet121-res224-all}       & $+0.530$ & $+0.151$ & $-0.082$ \\
\midrule
\multicolumn{4}{l}{\textit{Off-domain pretraining}} \\
\texttt{densenet121-res224-nih}       & $+0.373$ & $+0.313$ & $+0.007$ \\
\texttt{densenet121-res224-pc}        & $+0.435$ & $+0.313$ & $-0.050$ \\
\midrule
All 240 pairs                         & $\mathbf{+0.470}$ & $+0.196$ & $\mathbf{-0.056}$ \\
\bottomrule
\end{tabular}
\end{table}

% ============================================================
\section{CIFAR-10H: Theory-Confirming Negative Control}
\label{app:exp2-cifar10h}
% ============================================================

CIFAR-10H \citep{PetersonEtAl2019} evaluates seven pretrained CIFAR-10 \citep{Krizhevsky2009} classifiers \citep{ChenCIFARModels} against 50 human annotators across 10 one-vs-rest binary tasks (3{,}500 reader--task pairs).
Pretrained CIFAR-10 classifiers are highly accurate; their softmax probabilities rarely straddle a decision boundary for a given annotator's image subset.
Theorem~\ref{thm:facet-bridge} therefore predicts near-zero actionable boundary-regret mass everywhere.

\begin{table}[htbp]
\centering
\caption{CIFAR-10H negative control: mean Spearman $\bar\rho$ between log-loss gain and $\BRhat$ across $3{,}500$ reader--task pairs.}
\label{tab:cifar10h}
\small
\begin{tabular}{@{}lrrrr@{}}
\toprule
Reward & $\bar\rho$ (valid) & $\bar\rho$ (all) & $N$ valid & Top-decile $\BRhat=0$ \\
\midrule
R1 (thr.\ $0.500$) & $+0.015$ & $+0.007$ & $1{,}668$ & $96.2\%$ \\
R2 (thr.\ $0.167$) & $+0.044$ & $+0.024$ & $1{,}905$ & $94.5\%$ \\
R3 (thr.\ $0.833$) & $-0.008$ & $-0.004$ & $1{,}632$ & $96.9\%$ \\
\bottomrule
\end{tabular}
\end{table}

The R1 dissociation signal ($\bar\rho=+0.015$) is ${\sim}32{\times}$ weaker than the primary CheXpert R1 result ($\bar\rho=+0.477$ across the $25$ in-domain reader--task pairs; Section~\ref{sec:exp2}) and ${\sim}30{\times}$ weaker than the $240$-pair robustness average ($+0.470$; Table~\ref{tab:robustness}).
The collapse is the predicted consequence of dense actionable mass being absent, not a methods failure.

% ============================================================
\section{Allocation Advantage: Robustness}
\label{app:exp3-robustness}
% ============================================================

Table~\ref{tab:exp3-robustness} extends the allocation experiment of Section~\ref{sec:exp3} to all six DenseNet checkpoints and all eight CheXpert radiologists ($240$ reader--task pairs).
Three patterns reproduce the main result.
Under R1 (dense actionable mass), $\BRhat$ dominates Margin at every budget ($p<0.0001$ for $q\geq 10\%$), confirming the advantage is not an artefact of the primary checkpoint.
Under R2 (sparse but reachable mass), $\BRhat$ now also beats Margin at $q\in\{5\%,10\%,20\%\}$ ($p<0.0001$), an effect not visible on the $25$-pair primary benchmark; the dissociation table above explains why: off-domain checkpoints concentrate more belief mass near the R2 facet at $1/6$, providing actionable mass that the in-domain checkpoint lacks.
Under R3 (no actionable mass), $\BRhat$ and Margin remain economically indistinguishable, with no systematic advantage in either direction (sign-test $p \geq 0.08$ at every budget; mean differences $\leq 0.0012$ at $q=5\%$ and $\leq 0.0004$ for $q \geq 10\%$), reproducing the no-mass prediction.

\begin{table}[htbp]
\centering
\caption{Allocation robustness across $6$ DenseNet checkpoints $\times$ $8$ readers $\times$ $5$ conditions $=$ $240$ reader--task pairs.
  Entries report mean held-out utility gain over no review.
  W/L/T gives paired wins/losses/ties of $\BRhat$ against Margin; bold marks the best deployable policy in each row (Oracle excluded).
}
\label{tab:exp3-robustness}
\scriptsize
\setlength{\tabcolsep}{3pt}
\begin{tabular}{@{}ll rrrrr@{\hspace{5pt}}|@{\hspace{5pt}}rr@{\hspace{5pt}}|@{\hspace{5pt}}r@{}}
\toprule
Reward & $q$ & $\BRhat$ & Margin & Entropy & Residual & L2D & Random & Oracle & W/L/T \\
\midrule
R1 & $5\%$  & $\mathbf{+0.0407}$ & $+0.0403$ & $+0.0403$ & $+0.0192$ & $+0.0101$ & $+0.0195$ & $+0.0513$ & $106/74/60$ \\
R1 & $10\%$ & $\mathbf{+0.0787}$ & $+0.0763$ & $+0.0763$ & $+0.0413$ & $+0.0247$ & $+0.0393$ & $+0.0996$ & $134/71/35$ \\
R1 & $20\%$ & $\mathbf{+0.1534}$ & $+0.1444$ & $+0.1444$ & $+0.0939$ & $+0.0647$ & $+0.0788$ & $+0.1931$ & $167/51/22$ \\
R1 & $50\%$ & $\mathbf{+0.3160}$ & $+0.2936$ & $+0.2936$ & $+0.2653$ & $+0.2267$ & $+0.1947$ & $+0.4132$ & $198/25/17$ \\
\midrule
R2 & $5\%$  & $\mathbf{+0.0429}$ & $+0.0350$ & $+0.0368$ & $+0.0186$ & $+0.0258$ & $+0.0215$ & $+0.0604$ & $112/25/103$ \\
R2 & $10\%$ & $\mathbf{+0.0826}$ & $+0.0665$ & $+0.0685$ & $+0.0385$ & $+0.0541$ & $+0.0426$ & $+0.1078$ & $121/23/96$ \\
R2 & $20\%$ & $\mathbf{+0.1557}$ & $+0.1252$ & $+0.1276$ & $+0.0870$ & $+0.1123$ & $+0.0849$ & $+0.2013$ & $132/22/86$ \\
R2 & $50\%$ & $\mathbf{+0.3213}$ & $+0.3052$ & $+0.2682$ & $+0.2430$ & $+0.2826$ & $+0.2079$ & $+0.4399$ & $110/98/32$ \\
\midrule
R3 & $5\%$  & $\mathbf{+0.0042}$ & $+0.0030$ & $+0.0001$ & $+0.0008$ & $+0.0025$ & $+0.0002$ & $+0.0115$ & $22/31/187$ \\
R3 & $10\%$ & $\mathbf{+0.0052}$ & $+0.0048$ & $+0.0001$ & $+0.0013$ & $+0.0043$ & $+0.0003$ & $+0.0120$ & $18/30/192$ \\
R3 & $20\%$ & $+0.0055$ & $\mathbf{+0.0056}$ & $+0.0001$ & $+0.0020$ & $+0.0049$ & $+0.0005$ & $+0.0120$ & $13/23/204$ \\
R3 & $50\%$ & $+0.0057$ & $\mathbf{+0.0057}$ & $+0.0003$ & $+0.0025$ & $+0.0050$ & $+0.0017$ & $+0.0120$ & $8/18/214$ \\
\bottomrule
\end{tabular}
\end{table}

% TODO (revision appendix): Report the configured-system action-disagreement diagnostic: R1 versus R2 changes model-only actions on 32.5% of cases and post-review actions on 18.9%. %
